\begin{acknowledgments}
\chapter{Ringraziamenti}
Questa tesi è il frutto di ricerche e di lavoro svolto a conclusione della laurea magistrale.\\
Desidero, per primo, ringraziare il mio relatore, il Professor Valerio Leone Sciabolazza che mi ha fatto scoprire la mia passione per la ricerca e per l’economia dello sviluppo. Fin dal primo anno mi ha accompagnata con disponibilità e consigli preziosi, di cui farò tesoro.\\

Negli ultimi tempi molti paragonano la vita a un libro o a una serie TV (nel mio caso sarei al capitolo, o alla stagione, 24), ma non è un’immagine che sento davvero mia. Io penso alle persone che mi accompagneranno sempre, nel bene e nel male, e a quanto ho imparato da ciascuna e non riesco a vederli come personaggi secondari. Più che a un libro già scritto, preferisco paragonare la mia vita a un quaderno su cui continuo a scrivere: un quaderno un po’ particolare, che cambia stile a ogni pagina. Siamo noi, in fondo, a dare forma alle nostre avventure. \\
Sicuramente la mia famiglia è una di quelle cose su cui ho scritto e scriverò di più. Grazie a mamma e papà, che mi supportano nonostante non capiscano sempre le mie decisioni e che sono stati un riparo quando ne avevo bisogno. Fin da piccola mi hanno dato tutti gli strumenti per ritagliarmi un posto del mondo. \\
Grazie ai miei fratelli Francesco e Claudia. Non avremo il rapporto che spesso si vede nelle serie tv, ma negli anni ci stiamo avvicinando sempre di più. Grazie per aver "combattuto" molte battaglie al posto mio, ma anche grazie per avermi lasciato la possibilità di vincerne altre. Grazie Clo, per avermi insegnato che non si deve sapere tutto sempre dell'altra persona, ma che anche il silenzio vale molto e sopratutto per avermi insegnato che se si ha un obiettivo lo si può portare avanti costruendosi la propria strada. Grazie Checco per essere stato sempre sincero. Mi hai insegnato che se un'idea è giusta si deve combattere per cercare di realizzarla. Grazie per tutto l'aiuto con il codice, LaTex e per essermi stato vicino nei momenti di crisi (e, durante la tesi, non sono stati pochi). Insomma, grazie per l’esempio che mi avete dato, ciascuno a modo suo, nonostante le nostre differenze. Vi voglio bene. \\

Grazie a Fami Ohana, o famiglia padovana per meglio intenderci, che anche da lontano sono sempre i primi a festeggiare qualsiasi traguardo. Grazie zia Mock e zio Giovanni che mi fate ridere come non mai quando stiamo insieme, grazie zio Alberto che riesce sempre a farmi tornare bambina e grazie Zio Andrea e zia Caterina per essere presenti. Grazie Nonna Pia, mi hai insegnato moltissimo in questi anni, mi hai insegnato tantissimo in questi anni; grazie per essere sempre disponibile quando ti chiamo e per i consigli che mi dai. Sarai una nonna un po’ “atipica”, ma sei piena di vita e mi hai mostrato che, se si tiene la mente allenata, si può fare qualsiasi cosa si voglia. Un ringraziamento speciale a Nonno Nildo, che è sempre nei miei pensieri e che mi manca moltissimo. \\

Oltre la famiglia padovana, voglio ringraziare i miei vicini di casa, la famiglia tiburtina. Sono dei vicini un po' particolari, ma su cui si può sempre contare per qualsiasi cosa. Grazie Luigi, perchè sai sempre come strapparmi un sorriso, a casa sarà sempre presente un caffè con il dolcificante. Grazie zia Giò, i tutto: per essere sempre disponibile, anche nelle emergenze all’ultimo; per cercare di capirmi e sostenermi, ma anche per dirmi quando non sei d’accordo con me. Grazie, perché so che, anche se sono lontana e non chiamo spesso, sai che ti voglio bene. Ho tantissimi ricordi del tempo passato insieme: sono custoditi nel mio cuore, anche quando può sembrare che me ne dimentichi. Un grazie speciale anche a Nonna Maria, l'unica che è riuscita a farmi mangiare le verdure (fritte e pastellate) e che ricorderò sempre con la sua collana di perle e il sorriso. \\

Come tutti sappiamo, la famiglia non è fatta solo di legami “di sangue”: la famiglia si sceglie. Io la mia l’ho scelta molto tempo fa, o meglio, non l’ho scelta davvero, ed è questa la cosa più bella.\\
C'è una persona che è stata una costante nella mia vita, che mi ha fatto conoscere libri, serie, musica ma anche vivere avventure insieme. La mia persona. Carola. Mi sei stata vicina in tutti i momenti difficili, e in quelli bellissimi. Grazie Caro per avermi insegnato tanto e, onestamente, non saprei nemmeno da dove iniziare. Sei quella persona che, anche se non ci sentiamo per un mese, quando ci rivediamo è come se non fosse passato un giorno. Mi hai sempre incoraggiata a sfidare le mie capacità e a credere in me stessa, senza farmi fermare dalla paura di fallire. Con te passerei un'intera giornata a non fare nulla, solo a stare insieme: la tua compagnia mi fa stare bene, mi ricarica. La cosa che amo di più è venire a cena da te e cantare/gossippare mentre cucini. Grazie perchè senza il tuo supporto non avrei fatto metà delle cose che ho fatto, e di certo non sarei arrivata fin qui. Insieme abbiamo vissuto tantissime imprese e non vedo l'ora di viverne altre. Con te posso essere me stessa, e questo basta a rendere semplice perfino il silenzio. Come ti avevo detto tempo fa "voglio passare i vent'anni con te". Stiamo \textit{quasi} a metà percorso e continuo a pensare che non potessi scegliere meglio.\\

Ci sono delle persone con cui poi sono cresciuta insieme, anche se abbiamo preso strade molto diverse. Tizi e Mils, grazie per essere sempre voi stesse. Grazie Tizi per la tua spontaneità e ironia e Mils per la tua determinazione e per essere sempre solare. \\

Oltre alle amicizie storiche ce ne sono alcune che ho fatto nel percorso universitario e che hanno reso quest'esperienza unica. La prima fra tutti Sab. Ti ho conosciuta nel post-Covid, dopo un esame, per caso: la casualità migliore che potesse capitare. Sei stata un punto fermo, una persona con cui potevo essere sempre diretta. Mi hai insegnato a non arrendermi, che la propria strada si costruisce passo dopo passo e che nessuno, se non noi stesse, può giudicare le nostre scelte. Sono fiera della persona che sei diventata e io so che nessuno ti può fermare. Grazie per essermi stata vicina e aver saputo sempre cosa dire. \\
Insieme a Sab quell'anno ho conosciuto Alessio e Samuele, due colleghi universitari che non potrebbero essere più diversi tra loro. Ale, grazie perché sei sempre diretto e perché so che, se uno dei due ha bisogno, l’altro si farebbe in quattro. Le Pasquette più divertenti della mia vita portano il tuo nome. Sarò sempre qui per quando avrai bisogno. Samu, sei un paradosso per me, sei la persona più organizzata che conosco ma contemporaneamente quella più disorganizzata: ogni tanto scopro cose diverse che fai e non smetti mai di sorprendermi. Grazie per essere sempre stato un buon amico (e anche una buona spalla).\\ Sono davvero felice di aver passato questi 5 anni con voi ragazzi. \\

Grazie a Marco e Giovanni, due amici conosciuti l'ultimo anno di triennale che mi hanno cambiato completamente la visione dell'università. Gio, grazie per gli aperitivi e i caffè in cui mi fai morire dal ridere. Marco, te che ti metti in discussione ogni tre per due. Sei una persona fantastica, genuina, di quelle rare. Te che hai letteralmente deciso di adottarmi quando mi ha visto persa per i corridoi dell'università. Grazie, senza di te avrei perso tantissimo. \\

Grazie a EPOS sub ITA, senza che nomino tutti. Le lezioni grazie a voi sono state molto più leggere. Ci sono dei nomi speciali però che vorrei ringraziare: grazie Lucri, che nella tua ansia mi tranquillizzavi. Grazie per avermi fatta sentire capita fin da subito pisi. Grazie Bi, la tua gentilezza e il tuo ottimismo sono stati i motori di questi miei due anni, insieme ai nostri audio cortissimissimi. Da quando ti conosco non smetto di dire "ma è un cucciolo" e lo trovo bellissimo. Grazie Ale, sei stato un compagno di scleri, di "te hai iniziato?" "si sono a capitolo 5, te?" "no..." "a vabbè ma ce la fai". Sei una persona sincera, sorridente, non smettere di esserlo. Grazie Nora, per incoraggiarmi sempre e per capirmi al primo sguardo. Grazie di tutto ragazzi.  \\
Grazie a Francesco, che mi ha supportato e creduto in me sempre, soprattutto quando volevo mollare. Ha saputo incoraggiarmi, aiutarmi e ha sopportato molte, forse quasi troppe, mie sclerate relative sia all'università che non. Sei stato una sorpresa e mi hai insegnato che tutto si risolve, sempre: non tutto è grave come lo percepiamo e non tutto va preso sul serio. Grazie per ripetermi sempre: "una cosa alla volta, un passo dopo l'altro" quando mi sento troppo piena di cose da fare. Mi hai fatto riscoprire l'amore per l'economia, quella però più quotidiana, non solo quelle teorie che studiamo sui libri. Grazie perché so che posso contare su di te sempre, sopratutto quando non credo in me. Sei stato un punto di riferimento senza il quale avrei sicuro perso la rotta, e già io non ho un buon senso dell'orientamento. Il tuo sostegno in questi anni, ma sopratutto durante la stesura della tesi, è stato fondamentale. Grazie. E ricorda "tantissimo".\\

L'università, però, non è solo studio. E chi mi conosce lo sa, non perdo mai l'opportunità di riempirmi di cose da fare, sennò poi come mi lamento? \\
Proprio in università ho incontrato un mondo a cui non avrei mai pensato di appartenere e che, da spettatrice, avevo sempre ammirato. Casualmente, infatti, un giorno scopro che c'era una certa associazione che organizzava il TEDx Universitario. a scelta più pazza che potessi fare: io a organizzare eventi, a parlare con le aziende per ottenere fondi e riconoscimenti, a cercare tirocini tra il pubblico… io, che mi agito a scrivere un’email al professore.\\
Invece sono entrata, ed è stato amore a prima vista, ma sopratutto ho conosciuto delle persone stupende. \\
Per prima, Sabrina, che mi ha accolto a casa. Sei stata la presenza più importante del mio percorso associativo: una spalla quando tremavo, una consigliera quando serviva lucidità. Ci siamo conosciute online e, pagina dopo pagina, ti sei fatta spazio nel mio cuore. Ogni volta che ci vediamo il tempo corre troppo in fretta: sono grata di averti incontrata e di aver avuto la possibilità di conoscerti. Io sarò qui, a tifare per te, sempre. E poi… hai visto? Ho preparato una presentazione! \\
Oltre Sabrina, un grazie va a tutto il team ERS: Cecilia, Nicoletta e Matteo (e anche a chi ci siamo persi per la via). Ragazzi, non ho più parole per dire quanto sono orgogliosa di noi, del nostro percorso e del rapporto che abbiamo costruito. A prescindere so che riuscirete a fare ancora meglio. Mi mancheranno le nostre call, le strategie pensate insieme e anche i dibattiti su come rispondere alle richieste insolite degli sponsor.\\
Un grande grazie va anche al consiglio direttivo che mi ha sostenuta e insegnato tantissimo, accompagnandomi nel mio percorso. In particolare, grazie Silvia per aver sempre creduto in me e avermi consigliato, anche tardi la notte. Grazie per il sostegno: sei stata una bussola nei momenti confusi e una spinta quando avevo paura di fermarmi. Sei una persona speciale, che sa sempre cosa dire e ha una visione della vita stupenda, anche se diversa dalla mia. Con te mi sono sentita capita e al sicuro. Sono contenta di poterti considerare amica \\
Il TEDxSapienzaU non sarebbe lo stesso senza di voi. \\

A chi mi ha accompagnata fin qui: questa pagina porta anche la vostra firma.\\
Questa non è la fine del mio racconto, è una virgola. Il viaggio continua. Una cosa alla volta, un passo dopo l’altro: il mio quaderno resta aperto. Grazie
\end{acknowledgments}